% Generated by scripts/assemble.py; corrections belong in source/corrections.json.
\examhead{59}{天津大学 2025 年数学分析试卷(回忆版)}{58}
% p059-b20; source PDF page 59
\blockanchor{p059-b20}
\question{1}证明: $\int_{1}^{+\infty} \frac{\sin \pi x}{x} \, \mathrm{d}x < 0$.\par
% p059-b21; source PDF page 59
\blockanchor{p059-b21}
\question{2}$f(x)$ 是定义在 $[a, b]$ 上的函数, 且每一个点处极限都存在. 证明: $f(x)$ 在 $[a, b]$ 上有界.\par
% p059-b22; source PDF page 59
\par\Needspace{6\baselineskip}
\blockanchor{p059-b22}
\question{3}设 $f(x)$ 在 $[0, \pi]$ 上有连续可微, 且 $\int_0^\pi f(x) \mathrm{d}x = 0$. 证明:\par
\nopagebreak[4]
% p059-b23; source PDF page 59
\blockanchor{p059-b23}
\displayformula{\int_ {0} ^ {\pi} f ^ {2} (x) \mathrm{d} x \leq \int_ {0} ^ {\pi} [ f ^ {\prime} (x) ] ^ {2} \mathrm{d} x.}
% p059-b24; source PDF page 59
\par\Needspace{6\baselineskip}
\blockanchor{p059-b24}
\question{4}已知 $f(x)$ 在 $[0,1]$ 上二阶可导，且 $f(0) = 0, f(1) = 3, \min_{x \in [0,1]} f(x) = -1$. 证明：存在 $c \in (0,1)$，使得\par
\nopagebreak[4]
% p059-b25; source PDF page 59
\blockanchor{p059-b25}
\displayformula{f ^ {\prime \prime} (c) \geq 18.}
% p060-b00; source PDF page 60
\blockanchor{p060-b00}
\question{5}设 $x, y > 0$, 证明: $\frac{1}{4} (x^2 + y^2) \leq e^{x + 2y - 2}$.\par
% p060-b01; source PDF page 60
\blockanchor{p060-b01}
\question{6}设 $a_{n}, b_{n} \to 0 (n \to \infty)$, 且 $b_{1} > b_{2} > b_{3} > \dots$, 若 $\lim_{n \to \infty} \frac{a_{n} - a_{n-1}}{b_{n} - b_{n-1}} = A$, 证明: $\lim_{n \to \infty} \frac{a_{n}}{b_{n}} = A$.\par
% p060-b02; source PDF page 60
\blockanchor{p060-b02}
\question{7}设 $f(x)$ 是 $[0, +\infty)$ 上的可导函数, 且 $\lim_{x \to +\infty} f'(x) = A$ ($A$ 为有限数或 $+\infty$), 证明: $f(x)$ 在 $[0, +\infty)$ 上一致连续的充要条件是 $A$ 为有限数.\par
% p060-b03; source PDF page 60
\blockanchor{p060-b03}
\question{8}设 $\pmb{x} \in \mathbb{R}^3$, 定义范数 $N(\pmb{x}) : \mathbb{R}^3 \to \mathbb{R}$ 若满足\par
% p060-b04; source PDF page 60
\blockanchor{p060-b04}
(1) $N(\pmb{x}) \geq 0$, 且仅当 $\pmb{x} = \pmb{0}$ 时, 有 $N(\pmb{x}) = 0$.\par
% p060-b05; source PDF page 60
\blockanchor{p060-b05}
(2) $N(\pmb {x} + \pmb{y})\leq N(\pmb {x}) + N(\pmb {y}),\forall \pmb {x},\pmb {y}\in \mathbb{R}^{3}.$\par
% p060-b06; source PDF page 60
\blockanchor{p060-b06}
(3) $N(k\pmb {x}) = |k|N(\pmb {x}),\forall \pmb {x}\in \mathbb{R}^{3},k\in \mathbb{R}.$\par
% p060-b07; source PDF page 60
\blockanchor{p060-b07}
上述空间被称为欧氏空间. 内积定义的范数 $\| x\| = \sqrt{x\cdot x}$.\par
% p060-b08; source PDF page 60
\blockanchor{p060-b08}
(1) 证明: $N(\pmb{x})$ 在 $B = \{\pmb{x} \in \mathbb{R}^3 \mid \| \pmb{x} \| \leq 1\}$ 上有界.\par
% p060-b09; source PDF page 60
\blockanchor{p060-b09}
(2) $N(\pmb{x})$ 在 $\mathbb{R}^3$ 上连续.\par
% p060-b10; source PDF page 60
\par\Needspace{6\baselineskip}
\blockanchor{p060-b10}
(3) 对任意的 $x \in \mathbb{R}^3$, 记 $r = \| x \|$, 存在 $a, b \in \mathbb{R}$, 使得\par
\nopagebreak[4]
% p060-b11; source PDF page 60
\blockanchor{p060-b11}
\displayformula{a r \leq N (\boldsymbol {x}) \leq b r.}
% p060-b12; source PDF page 60
\blockanchor{p060-b12}
\question{9}记 H 为黑塞矩阵.\par
% p060-b13; source PDF page 60
\blockanchor{p060-b13}
(1) 在点 $(a, b)$ 处, 已知 $f_{x}(a, b) = f_{y}(a, b) = 0$, $H = \left( \begin{array}{ll} f_{xx}(x, y) & f_{xy}(x, y) \\ f_{yx}(x, y) & f_{yy}(x, y) \end{array} \right)$, 且 $H$ 在 $(a, b)$ 处正定, 证明: $(a, b)$ 为 $f$ 的严格极小值点.\par
% p060-b14; source PDF page 60
\blockanchor{p060-b14}
(2) 若 $H$ 在 $\mathbb{R}^2$ 上处处正定, 则 $f$ 在 $\mathbb{R}^2$ 上至多有一个极值点.\par
% p060-b15; source PDF page 60
\blockanchor{p060-b15}
\question{10}设 $B$ 为单位球, $f$ 在 $B$ 上有定义, $\pmb{v} = (v_{1}, v_{2}, v_{3})$, $\nabla \cdot \pmb{v} = 0$. 在 $B$ 上有 $\pmb{v} = 0$.\par
% p060-b16; source PDF page 60
\blockanchor{p060-b16}
(1) 求证: $\int_{B} \nabla f \cdot \boldsymbol{v} \, \mathrm{d}x = 0$.\par
% p060-b17; source PDF page 60
\blockanchor{p060-b17}
(2) 求证:  $\int_{B} v_{1} dx = 0.$\par
